\section{Foundations and the cases \texorpdfstring{$4\leq n\leq9$}{4 <= n <= 9}}
\label{sec:foundations-small}

This section fixes the conventions used throughout the paper, derives the
universal counting rows, and proves all values up to nine.  A line is used as
a bookkeeping carrier, but is never counted as a circle.

\subsection{The problem and the block language}

\begin{definition}[The V1 problem]\label{def:v1-admissible}
Let $P\subset\R^2$ be a set of $n\geq4$ distinct points.  We call $P$
\emph{admissible} if it is neither collinear nor contained in one proper
Euclidean circle.  A \emph{determined circle} is a proper circle containing a
noncollinear triple of $P$.  Write $C(P)$ for the number of distinct
determined circles and
\[
  f(n):=\min\{C(P): |P|=n\text{ and }P\text{ is admissible}\}.
\]
\end{definition}

For a line or a proper circle $X$, its \emph{full trace} is $P\cap X$.
A maximal generalized $s$-block is a full trace of size $s\geq3$.  We write
\[
 L_s:=\#\{\text{$s$-point line blocks}\},\qquad
 C_s:=\#\{\text{$s$-point circle blocks}\},\qquad
 B_s:=L_s+C_s.
\]
Thus
\begin{equation}\label{eq:circle-objective}
 C(P)=\sum_{s\geq3}C_s.
\end{equation}
For $p\in P$, let $\ell_s(p)$, $c_s(p)$, and
$d_s(p)=\ell_s(p)+c_s(p)$ count, respectively, the line, circle, and all
$s$-blocks through $p$.

\begin{lemma}[Unique ownership]\label{lem:triple-ownership}
Every triple of $P$ belongs to exactly one maximal generalized block.
Consequently
\begin{align}
 \mathsf T
   &:=\sum_{s\geq3}\binom{s}{3}B_s=\binom n3,
       \label{eq:row-T}\\
 \sum_{p\in P}d_s(p)&=sB_s,
 &\sum_{p\in P}\ell_s(p)&=sL_s.                  \label{eq:block-incidence}
\end{align}
If $L_2$ denotes the number of ordinary, two-point full line traces, then
\begin{equation}\label{eq:line-pair-partition}
 L_2+\sum_{s\geq3}\binom{s}{2}L_s=\binom n2.
\end{equation}
\end{lemma}

\begin{proof}
A collinear triple has its unique carrier line; a noncollinear triple has its
unique circumcircle.  Passing to the full trace gives a unique maximal owner.
Equation~\eqref{eq:row-T} counts triples by that owner, and
\eqref{eq:block-incidence} counts point--block incidences.  Finally, every
pair of distinct points has one carrier line, which proves
\eqref{eq:line-pair-partition}.
\end{proof}

\begin{lemma}[The eight-point Orchard bound]\label{lem:orchard-eight}
Let \(Q\) be a set of eight noncollinear points of the real projective plane,
with no four collinear.  If \(t_i\) denotes the number of lines containing
exactly \(i\) points of \(Q\), then
\[
 t_3\leq 7.
\]
\end{lemma}

\begin{proof}
Pair counting gives \(t_2+3t_3=28\).  Melchior's inequality excludes
\(t_3\geq9\), so it remains only to rule out \(t_3=8\).  In that case
\(t_2=4\).  At a point \(v\in Q\), let \(r_v\) and \(s_v\) be the numbers
of three-point and ordinary lines through \(v\).  Then
\[
 7=2r_v+s_v.
\]
Thus \(s_v\) is a positive odd integer, while
\(\sum_v s_v=2t_2=8\).  Consequently \(s_v=1\) and \(r_v=3\) for every
\(v\).  The four ordinary lines are therefore a perfect matching; denote
its pairs by
\[
 A=\{A_0,A_1\},\quad B=\{B_0,B_1\},\quad
 C=\{C_0,C_1\},\quad D=\{D_0,D_1\}.
\]
Every three-point line uses three different pairs.  Since the eight such
lines account for all \(24\) cross-pair incidences, exactly two omit each
of \(A,B,C,D\).  Explicitly, if \(n_A,n_B,n_C,n_D\) are the four
omission counts, then the four edges between, say, \(A\) and \(B\) give
\(n_C+n_D=4\); the six analogous equations force
\(n_A=n_B=n_C=n_D=2\).

We first determine the resulting incidence structure.  Two triples on the
same three pairs cannot differ in only one entry, because then their two
common points would determine both lines.  They also cannot differ in
exactly two entries.  Indeed, after relabelling suppose the two \(ABC\)
triples are
\[
 A_0B_0C_0,\qquad A_1B_1C_0.
\]
The unused cross edges force the remaining triples to have the form
\[
\begin{array}{ll}
 A_0B_1D_x,&A_1B_0D_y,\\
 A_0C_1D_z,&A_1C_1D_{1-z},\\
 B_0C_1D_w,&B_1C_1D_{1-w}.
\end{array}
\]
The \(AD\), \(BD\), and \(CD\) edges successively give
\(y=z=1-x\) and \(w=x\), after which one of the two \(C_1D_i\) edges is
repeated.  Hence the two triples on any fixed three pairs are complementary.
After fixing the \(ABC\) pair as \(000,111\), the unused \(AB,AC,BC\)
edges are the three cross-matchings.  A flip of the \(D\)-labels fixes one
\(ABD\) triple, and the unused \(AD\) and \(BD\) edges then force all
remaining bits.  Up to relabelling, the unique incidence list is
\begin{equation}\label{eq:orchard-eight-incidence}
\begin{split}
&A_0B_0C_0,\quad A_1B_1C_1,\quad
 A_0B_1D_0,\quad A_1B_0D_1,\\
&A_0C_1D_1,\quad A_1C_0D_0,\quad
 B_0C_1D_0,\quad B_1C_0D_1.
\end{split}
\end{equation}

Choose the projective frame
\[
 A_0=(1,0,0),\quad A_1=(0,1,0),\quad
 B_0=(0,0,1),\quad B_1=(1,1,1).
\]
The first four incidences in \eqref{eq:orchard-eight-incidence} allow the
normalization
\[
 C_0=(1,0,u),\quad C_1=(1,v,1),\quad
 D_0=(w,1,1),\quad D_1=(0,1,z),
\]
where the relevant parameters are nonzero.  The next incidences give
\[
 vz=uw=vw=1.
\]
Thus \(u=v\) and \(w=z=1/u\).  The last incidence,
\(B_1C_0D_1\), now gives
\[
 u^2-u+1=0,
\]
which has no real root.  Hence \(t_3=8\) is impossible.
\end{proof}

\begin{lemma}[Radical-axis cross-block bound]
\label{lem:radical-axis-cross-block}
Let \(A\) and \(B\) be marked points on two distinct real circles, after
deleting any marked point common to the circles.  A \emph{cross-block} is a
line or circle containing two points of \(A\) and two points of \(B\).
If \((|A|,|B|)=(5,4)\), there are at most \(12\) cross-blocks.  If
\((|A|,|B|)=(4,4)\), there are at most \(10\).
\end{lemma}

\begin{proof}
For each chord of the first circle and each chord of the second, the four
endpoints are concyclic precisely when the two supporting lines meet on the
radical axis.  A coincident pair of supporting lines gives a line
cross-block and is assigned to its intersection with the same radical
axis.  (The common supporting line cannot itself be the radical axis,
because the marked common points were deleted.)

At a point \(q\) of the radical axis, let \(a_q\) and \(b_q\) be the
numbers of selected chords from \(A\) and \(B\) through \(q\).  Chords
through a fixed external point form a matching, so
\[
 a_q\leq\floor{|A|/2},\qquad b_q\leq\floor{|B|/2}.
\]
Moreover, \(\sum_q a_q=\binom{|A|}{2}\),
\(\sum_q b_q=\binom{|B|}{2}\), and the number of cross-blocks is
\[
 M=\sum_q a_qb_q.
\]
For \(5+4\), this gives \(M\leq2\sum_qb_q=12\).

For \(4+4\), both multiplicity sums are \(6\), and \(M\leq12\).  Equality
\(M=11\) is impossible: from
\[
 12-M=\sum_q a_q(2-b_q)=1
\]
there would be one point with \(a_q=b_q=1\), while every other point with
\(a_q>0\) had \(b_q=2\).  If \(r\) and \(s\) count the remaining points
with \(a_q=1\) and \(a_q=2\), respectively, then
\[
 1+r+2s=6,\qquad 1+2r+2s\leq6,
\]
forcing \(r=0\) and then \(1+2s=6\), a contradiction.  Equality \(M=12\)
would force exactly three points \(q\), with
\(a_q=b_q=2\) at each.  These are the three diagonal points of each of the
two complete quadrangles, yet all three lie on the radical axis.  The
diagonal points of a nondegenerate complete quadrangle are noncollinear
(in the standard frame their determinant is \(-2\)).  Thus \(M\leq10\).
\end{proof}

\subsection{Sharp upper constructions}

\begin{construction}[A circle and its centre]\label{con:circle-centre}
Put $r=n-1$ distinct points on a proper circle $\Gamma$ and add its centre
$o$.  Choose exactly $\floor{r/2}$ antipodal boundary pairs; if $r$ is odd,
choose the remaining boundary point without its antipode.  The boundary
triples determine only $\Gamma$.  A triple consisting of $o$ and a boundary
pair is collinear precisely for an antipodal pair.  Every other boundary
pair gives a different proper circle, since a circle different from
$\Gamma$ meets $\Gamma$ in at most two points.  Thus this admissible set
determines exactly
\begin{equation}\label{eq:circle-centre-count}
 F(n):=1+\binom{n-1}{2}-\floor{\frac{n-1}{2}}
\end{equation}
 circles.  In particular it gives the sharp upper targets $3,5$ for
 $n=4,5$ and $25$ for $n=9$.
\end{construction}

 The three exceptional upper bounds have equally explicit fixed certificates.
For points $p_i=(x_i,y_i)$, write
\begin{equation}\label{eq:concyclic-determinant}
 \Delta(i,j,k,l):=
 \det\begin{pmatrix}
 x_i^2+y_i^2&x_i&y_i&1\\
 x_j^2+y_j^2&x_j&y_j&1\\
 x_k^2+y_k^2&x_k&y_k&1\\
 x_l^2+y_l^2&x_l&y_l&1
 \end{pmatrix}.
\end{equation}
For a noncollinear quadruple, $\Delta=0$ is equivalent to concyclicity.

\begin{construction}[Six rational points and eight circles]
\label{con:six-upper}
Take
\[
\begin{array}{c|cccccc}
i&0&1&2&3&4&5\\ \hline
p_i&(0,0)&(4,0)&(1,0)&(1,3)&
 (\tfrac25,\tfrac65)&(2,2).
\end{array}
\]
The collinear triples are exactly
\[
 012,\qquad034,\qquad135.
\]
Exact substitution in $x^2+y^2+Dx+Ey+F=0$ gives the following eight
circle traces and coefficient triples $(D,E,F)$:
\[
\begin{array}{c|c}
013&(-4,-2,0)\\
0145&(-4,0,0)\\
0235&(-1,-3,0)\\
024&(-1,-1,0)\\
1234&(-5,-3,4)\\
125&(-5,-1,4)\\
245&(-3,-2,2)\\
345&(-2,-4,4).
\end{array}
\]
The three four-circles own four triples each and the other five circles own
one triple each.  Together with the three collinear triples this accounts for
all $\binom63=20$ triples.  Thus the configuration is admissible and
determines exactly $8$ circles.
\end{construction}

\begin{construction}[Seven points and eleven circles]
\label{con:seven-upper}
Let $u=\sqrt3$ and take
\[
\begin{array}{c|ccccccc}
i&0&1&2&3&4&5&6\\ \hline
p_i&(-1,0)&(1,0)&(0,u)&(0,0)&
(-\tfrac12,\tfrac u2)&(\tfrac12,\tfrac u2)&(0,\tfrac u3).
\end{array}
\]
The collinear triples are exactly
\[
 013,\ 024,\ 056,\ 125,\ 146,\ 236,
\]
where a string denotes the corresponding set of indices.  Direct expansion
of \eqref{eq:concyclic-determinant}, using the dihedral symmetry of the
equilateral triangle, gives exactly the six proper four-circles
\[
 0145,\ 0235,\ 0346,\ 1234,\ 1356,\ 2456.
\]
Their $24$ triples are disjoint.  The five remaining noncollinear triples
are
\[
 012,\ 016,\ 026,\ 126,\ 345.
\]
The determinant list also shows that they have five different owners and
that no five-point circle occurs.  Hence the configuration determines
$6+5=11$ circles.
\end{construction}

\begin{construction}[Eight rational points and seventeen circles]
\label{con:eight-upper}
Take
\[
\begin{array}{c|cccccccc}
i&0&1&2&3&4&5&6&7\\ \hline
p_i&(0,0)&(1,0)&(\tfrac13,0)&(\tfrac23,0)&
(\tfrac23,\tfrac13)&(\tfrac25,\tfrac15)&
(\tfrac13,\tfrac13)&(\tfrac35,\tfrac15).
\end{array}
\]
Its collinear triples are exactly
\[
 012,\ 013,\ 023,\ 123,\ 045,\ 167.
\]
Exact rational evaluation of \eqref{eq:concyclic-determinant} gives exactly
the eleven proper four-circles
\[
\begin{gathered}
0146,\ 0157,\ 0247,\ 0256,\ 0367,\ 1245,\\
1347,\ 1356,\ 2346,\ 2357,\ 4567.
\end{gathered}
\]
No five-point subset has all its quadruples on this list.  The six remaining
noncollinear triples are
\[
 034,\ 035,\ 126,\ 127,\ 267,\ 345.
\]
Thus the eleven four-circles own $44$ triples and the last six triples have
six distinct owners, for a total of $17$ circles.
\end{construction}

\subsection{The elementary cases \texorpdfstring{$n=4,5$}{n=4,5}}

\begin{proposition}\label{prop:n4-n5}
One has $f(4)=3$ and $f(5)=5$.
\end{proposition}

\begin{proof}
Let first $|P|=4$.  If three points lie on a line, the fourth point together
with each of the three pairs on that line gives a proper circle; the three
circles are distinct.  If no triple is collinear, all four triples are
noncollinear, and two of them could share a circle only if all four points
were concyclic.  Hence $C(P)\geq3$.  Construction
\ref{con:circle-centre} attains three.

Now let $|P|=5$.  A four-point line and its outsider give six distinct
circles, so assume that no line has four points.  Distinct three-point lines
use disjoint pairs; therefore $3L_3\leq\binom52$ and $L_3\leq3$.  If no
four-point circle occurs, every noncollinear triple has a different owner,
and
\[
 C(P)=\binom53-L_3\geq7.
\]
If a four-point circle $\Gamma$ occurs, it is unique.  Every three-point line
contains the outsider and a chord of $\Gamma$, and these chords form a
matching, so there are at most two such lines.  At least eight triples are
noncollinear; four lie on $\Gamma$, while the other four have mutually
different owners, since a repeated owner would be a second four-point
circle.  Thus $C(P)\geq1+4=5$, attained by
\cref{con:circle-centre}.
\end{proof}


We repeatedly use, without further comment, that two distinct proper circles
meet in at most two points, two distinct lines in at most one point, and a
line meets a proper circle in at most two points.

\subsection{Classical incidence inputs and their exact scope}
\label{subsec:classical-inputs}

For later reference we record every external incidence theorem used in the
proof.  The hypotheses below are part of the statements; in particular, none
of these results assumes general position.

\begin{theorem}[Melchior]\label{thm:melchior}
Let $\mathcal A$ be a finite arrangement of distinct real projective lines,
not all concurrent, and let $t_r(\mathcal A)$ count vertices incident with
exactly $r$ lines.  Then
\begin{equation}\label{eq:melchior}
 t_2\geq3+\sum_{r\geq4}(r-3)t_r.
\end{equation}
Equivalently, the same inequality holds for the multiplicities of lines
determined by a finite set of distinct noncollinear real projective points.
All projective vertices, including those at infinity, are counted.
\end{theorem}

This is the point-dual form of Melchior's theorem~\cite{Melchior1940}.  We do
not import a classification of its equality cases.

\begin{theorem}[Kelly--Moser]\label{thm:kelly-moser}
If $Q$ is a finite set of $q$ distinct noncollinear real points and $t_2$ is
the number of its ordinary determined lines, then
\begin{equation}\label{eq:kelly-moser}
 3q\leq7t_2,
 \qquad
 t_2\geq\ceil{\frac{3q}{7}}.
\end{equation}
The inequality is non-strict; no equality classification is assumed.
\end{theorem}

This is Theorem~3.6 of Kelly and Moser~\cite{KellyMoser1958}.

\begin{theorem}[Langer--de Zeeuw form]\label{thm:langer}
Let $Q$ be a finite set of $N$ distinct affine points over $\mathbb C$.
Assume that every affine line contains at most $2N/3$ points of $Q$.  If
$t_r$ counts the spanned lines containing exactly $r$ points, then
\begin{equation}\label{eq:langer}
 \sum_{r\geq2}r t_r\geq\frac{N(N+3)}{3}.
\end{equation}
In particular the result applies to real point sets.  Both the occupancy cap
and the conclusion are non-strict.
\end{theorem}

We use the formulation of de Zeeuw~\cite{DeZeeuw2018}, derived from
Langer's orbifold inequality~\cite{Langer2003}.

\begin{theorem}[Lenchner]\label{thm:lenchner}
Let $\mathcal A$ be an affine arrangement of $q$ distinct real lines, neither
all parallel nor all concurrent.  Apart from one specified exceptional
six-line arrangement, the number $o_{\mathrm{fin}}(\mathcal A)$ of finite
ordinary intersections satisfies
\begin{equation}\label{eq:lenchner}
 7o_{\mathrm{fin}}(\mathcal A)\geq2q-3.
\end{equation}
Consequently \eqref{eq:lenchner} is available under the simpler guard
$q\neq6$.
\end{theorem}

The word \emph{finite} is essential.  This is the corollary of Lenchner's
Theorem~8~\cite{Lenchner2007} used later in the paper.

\begin{theorem}[Ungar, six-point consequence]\label{thm:ungar-six}
Six distinct noncollinear real affine points determine at least six
directions, where a direction is a parallelism class of determined lines.
\end{theorem}

This is the only consequence of Ungar's direction theorem that we use
\cite{Ungar1982}.

\begin{theorem}[Hall]\label{thm:hall}
Let $G$ be a finite bipartite graph with parts $X,Y$.  If
$|S|\leq|N(S)|$ for every $S\subseteq X$, then $G$ has a matching saturating
$X$.
\end{theorem}

We use Hall's theorem~\cite{Hall1935} only after the relevant neighbourhood
inequalities have been established geometrically.

Inversion, projective duality, directed Menelaus, and power of a point are
not treated as external incidence black boxes.  The inversion statements
needed for the universal rows are proved next; later geometric uses are
proved at their point of application.

\subsection{Inversion and the universal rows}

Fix $p\in P$ and $\rho>0$.  On $\R^2\setminus\{p\}$ put
\[
 I_{p,\rho}(x)=p+\frac{\rho}{\|x-p\|^2}(x-p).
\]

\begin{lemma}[Inversion dictionary]\label{lem:inversion-dictionary}
The map $I_{p,\rho}$ is an injective involution on
$\R^2\setminus\{p\}$.  A line through $p$ maps to itself, while a proper
circle through $p$ maps to a line avoiding $p$, and conversely.  Hence every
maximal $s$-block through $p$ becomes a full line containing exactly $s-1$
points of the inverted set
\[
 Q_p:=I_{p,\rho}(P\setminus\{p\}).
\]
The set $Q_p$ is noncollinear.  It remains noncollinear after the inversion
centre $o=p$ is added.
\end{lemma}

\begin{proof}
The formula immediately gives $I_{p,\rho}^2=\mathrm{id}$.  Translating
$p$ to the origin and substituting the formula into a linear or quadratic
equation gives the stated line--circle dictionary.  Full traces are
preserved because the map is injective.

If $Q_p$ lay on a line through $o$, then all original points would lie on
that line.  If it lay on a line avoiding $o$, the inverse carrier would be a
proper circle through $p$ containing all of $P$.  Both possibilities violate
admissibility.  Adding $o$ cannot make a noncollinear set collinear.
\end{proof}

\begin{corollary}[The Lenchner pivot bridge]\label{cor:pivot-lenchner}
Let \(p\in P\), put \(q=n-1\), and assume \(q\neq6\).  Then at least
\[
 \ceil{\frac{2q-3}{7}}
\]
proper three-point circles of \(P\) pass through \(p\).  Equivalently,
\[
 c_3(p)\geq\ceil{\frac{2n-5}{7}}.
\]
\end{corollary}

\begin{proof}
Invert \(P\setminus\{p\}\) at \(p\), obtaining the \(q\) distinct,
noncollinear image points \(Q_p\) from
\cref{lem:inversion-dictionary}.  Work projectively, dualize those points,
and declare the dual \(p^*\) of the inversion centre to be the line at
infinity.  This produces \(q\) distinct real affine lines.

The dual lines are not all parallel: otherwise the points of \(Q_p\)
would be collinear on a line through \(p\).  They are not all concurrent
at a finite point: otherwise \(Q_p\) would be collinear on a line avoiding
\(p\).  Thus the arrangement satisfies the geometric hypotheses of
\cref{thm:lenchner}; the guard \(q\neq6\) also removes its sole exceptional
cardinality.  It consequently has at least
\(\ceil{(2q-3)/7}\) finite ordinary intersections.

Under projective duality, such an intersection is an ordinary determined
line of \(Q_p\).  Finiteness says exactly that this line avoids \(p\):
ordinary intersections on \(p^*\) correspond instead to image lines
through \(p\), and are deliberately not counted.  By inversion, an ordinary
image line avoiding \(p\) is a proper circle whose full trace on \(P\)
consists of \(p\) and the two preimages.  Distinct image lines give distinct
circles, proving the claim.
\end{proof}

The inversion dictionary verifies every noncollinearity hypothesis in the next
two applications of \cref{thm:melchior}.

\begin{proposition}[Universal Melchior rows]\label{prop:universal-melchior-rows}
For every admissible $n$-point set and every $p\in P$, define
\begin{equation}\label{eq:pivot-slack}
 \sigma_p:=d_3(p)-3-\sum_{s\geq5}(s-4)d_s(p).
\end{equation}
Then $\sigma_p\geq0$, and
\begin{align}
 \mathsf P
   &:=\sum_{s\geq3}s(4-s)B_s\geq3n,
      \label{eq:row-P}\\
 \sum_{p\in P}\sigma_p&=\mathsf P-3n.             \label{eq:sum-pivot-slack}
\end{align}
If $\ell_2(p)$ counts the ordinary original lines through $p$, then
\begin{align}
 \ell_2(p)+\sum_{s\geq3}(s-1)\ell_s(p)&=n-1,
      \label{eq:pivot-arms}\\
 \sum_{s\geq3}s\ell_s(p)&\leq n-1+\sigma_p.
      \label{eq:added-centre-local}
\end{align}
After summing \eqref{eq:added-centre-local}, one obtains
\begin{equation}\label{eq:row-D}
 \mathsf D
 :=\sum_{s\geq3}s(s-4)C_s
   +\sum_{s\geq3}2s(s-2)L_s
 \leq n(n-4).
\end{equation}
Finally, direct application of Melchior to $P$ gives
\begin{equation}\label{eq:global-line-row}
 3L_3+\sum_{s\geq4}
   \left(\binom{s}{2}+s-3\right)L_s
 \leq\binom n2-3.
\end{equation}
\end{proposition}

\begin{proof}
Apply \cref{thm:melchior} to the $n-1$ inverted points.  By
\cref{lem:inversion-dictionary}, its $(s-1)$-point lines are exactly the
$s$-blocks through $p$.  Melchior is therefore precisely
$\sigma_p\geq0$.  Summing over $p$ and using
\eqref{eq:block-incidence} gives \eqref{eq:row-P} and
\eqref{eq:sum-pivot-slack}.

The original lines through $p$ partition $P\setminus\{p\}$ into arms,
which proves \eqref{eq:pivot-arms}.  Add the inversion centre $o$.  The
line multiplicities of the resulting $n$-point set are exactly
\[
 t_2=c_3(p)+\ell_2(p),\qquad
 t_i=c_{i+1}(p)+\ell_i(p)\quad(i\geq3).
\]
The slack in Melchior is then
\[
 \sigma_p+\ell_2(p)-\sum_{s\geq3}\ell_s(p)\geq0.
\]
Using \eqref{eq:pivot-arms} yields
\eqref{eq:added-centre-local}.  Sum it over $p$, use
\eqref{eq:block-incidence} and \eqref{eq:sum-pivot-slack}, and subtract
$\mathsf P$ from the resulting inequality.  The circle coefficient is
$-s(4-s)=s(s-4)$ and the line coefficient is
$s^2-s(4-s)=2s(s-2)$, giving \eqref{eq:row-D}.

For the last assertion, substitute \eqref{eq:line-pair-partition} in
\eqref{eq:melchior} applied directly to $P$.
\end{proof}

\begin{proposition}[Conditional Langer row]\label{prop:universal-langer-row}
Suppose that, for every pivot $p$, every line determined by $Q_p$ contains
at most $2(n-1)/3$ image points.  Then
\begin{equation}\label{eq:row-L}
 \mathsf L:=\sum_{s\geq3}s(s-1)B_s
 \geq\frac{n(n-1)(n+2)}{3}.
\end{equation}
\end{proposition}

\begin{proof}
Apply \cref{thm:langer} to the $N=n-1$ image points at each pivot.  Its
incidence count is $\sum_{s\geq3}(s-1)d_s(p)$.  Summing over $p$ and using
\eqref{eq:block-incidence} proves \eqref{eq:row-L}.  Notice that the
occupancy hypothesis must be verified before this row is used.
\end{proof}

Equations \eqref{eq:row-T}, \eqref{eq:row-P}, \eqref{eq:row-D}, and,
under its stated cap, \eqref{eq:row-L}, are the four universal rows imported
by the later sections.

\subsection{Two pencil bounds}

\begin{lemma}[Rich-line pencil]\label{lem:rich-line-pencil}
If a line contains $m\geq3$ points of an admissible $n$-point set and
$q=n-m\geq1$ points lie off it, then
\begin{equation}\label{eq:rich-line-pencil}
 C(P)\geq q\binom m2-\binom q2\floor{\frac m2}.
\end{equation}
\end{lemma}

\begin{proof}
For a fixed outsider, its circles with pairs on the line are proper and
mutually distinct.  For two fixed outsiders, common circles use disjoint
pairs of line points: two pairs sharing a point would give two circles with
three common points.  Thus two outsider pencils overlap in at most
$\floor{m/2}$ circles.  The two-term Bonferroni inequality gives
\eqref{eq:rich-line-pencil}; higher intersections do not invalidate this
lower bound.
\end{proof}

\begin{lemma}[Rich-circle pencil]\label{lem:rich-circle-pencil}
If a proper circle $\Gamma$ contains $m\geq3$ selected points and
$q=n-m\geq1$ points lie outside it, then
\begin{equation}\label{eq:rich-circle-pencil}
 C(P)\geq
 1+q\left(\binom m2-\floor{\frac m2}\right)
   -\binom q2\floor{\frac m2}.
\end{equation}
\end{lemma}

\begin{proof}
For an outsider $x$, the chords of $\Gamma$ passing through $x$ form a
matching on the $m$ marked points, so at most $\floor{m/2}$ marked pairs are
collinear with $x$.  Every other pair gives a distinct proper circle through
$x$.  As in the line case, two outsider pencils overlap in at most
$\floor{m/2}$ circles.  Add $\Gamma$ and apply Bonferroni.
\end{proof}

\subsection{Six points}

\begin{proposition}\label{prop:n6}
One has $f(6)=8$.
\end{proposition}

\begin{proof}
Construction~\ref{con:six-upper} gives the upper bound.  Suppose, for a
contradiction, that an admissible six-point set satisfies $C(P)\leq7$, and
let $k$ be the maximum number of selected points on a proper circle.

If $k=5$, \cref{lem:rich-circle-pencil} gives $C(P)\geq9$.  If $k=3$,
every noncollinear triple has a different circle owner.  Let a largest line
contain $m$ points and put $r=6-m$.  For $m\geq4$, a collinear triple not
contained in that line uses at most one of its points, and an outsider pair
can be collinear with at most one point on the line.  Thus the number $T_0$
of collinear triples satisfies
\[
 T_0\leq\binom m3+\binom r3+\binom r2.
\]
This leaves at least $10$ noncollinear triples for $m=5$ and at least $15$
for $m=4$.  If all lines have size at most three, their collinear triples use
disjoint sets of three pairs, so there are at most
$\floor{\binom62/3}=5$ of them.  Again at least $15$ circles remain.

It remains that $k=4$.  The rich-line bound is $10$ for both $m=4$ and
$m=5$, so every line has size at most three.  All maximal generalized blocks
therefore have size three or four, and every four-block is a circle.  Put
\[
 b=C_4,\qquad L=L_3.
\]
Triple ownership, the pivot Melchior row, the exact circle count, and the
global line row give, respectively,
\begin{align}
 B_3+4b&=20,                                      \label{eq:n6-triples}\\
 0\leq\sum_{p\in P}(d_3(p)-3)&=42-12b,           \label{eq:n6-pivot}\\
 C(P)&=20-3b-L,                                   \label{eq:n6-circles}\\
 L&\leq4.                                         \label{eq:n6-lines}
\end{align}
Hence $b\leq3$, while $C(P)\leq7$ forces
\begin{equation}\label{eq:n6-terminal-rows}
 b=3,\qquad L=4.
\end{equation}

At any $p\in P$, ownership of the ten triples containing $p$ reads
\[
 d_3(p)+3d_4(p)=10.
\]
The first Melchior row gives $d_3(p)\geq3$, and $d_3(p)\equiv1\pmod3$.
Since $\sum_p d_4(p)=4b=12$, necessarily
\begin{equation}\label{eq:n6-local-profile}
 d_3(p)=4,\qquad d_4(p)=2,\qquad \sigma_p=1
 \quad\text{for every }p.
\end{equation}
The added-centre inequality \eqref{eq:added-centre-local} becomes
\begin{equation}\label{eq:n6-kappa}
 \kappa_p:=6-3\ell_3(p)\geq0.
\end{equation}

Since $L=4$, summing \eqref{eq:n6-kappa} gives zero, so every point lies on
exactly two of the four three-point lines.  No three of those lines are
concurrent, and their six pairwise intersections are the selected points:
the line pattern is a complete quadrilateral.  Label its four lines
\[
 012,\qquad034,\qquad135,\qquad245.
\]
Apply a Euclidean similarity, which preserves lines and circles, so that
\begin{equation}\label{eq:n6-pasch-coordinates}
\begin{gathered}
p_0=(0,0),\quad p_1=(1,0),\quad p_3=(u,v),\qquad v\ne0,\\
p_2=(a,0),\quad p_4=b(u,v),\quad
p_5=t(1,0)+(1-t)(u,v),\qquad
t=\frac{a(b-1)}{b-a}.
\end{gathered}
\end{equation}
Here distinctness gives $a,b\notin\{0,1\}$.  Moreover $a\neq b$: if
$a=b$, then the line $p_2p_4$ is parallel to $p_1p_3$ (their direction
vectors differ by the nonzero scalar $a$), and the two lines are distinct
because $a\neq1$; they therefore cannot share the affine point $p_5$.
Put
$R=u^2+v^2>0$.

By \eqref{eq:n6-local-profile}, the complements of the three four-circles
are three disjoint pairs partitioning $P$.  Each complement meets every one
of the four collinear triples.  In a complete quadrilateral the only such
pairs are the three diagonal pairs, so the circle supports are forced to be
\[
 0145,\qquad0235,\qquad1234.
\]
Substitution in \eqref{eq:concyclic-determinant} gives, in the same order,
\[
\begin{split}
 \Delta(0,1,4,5)
  &=\frac{vb^2(a-1)(b-1)}{(a-b)^2}(-Rb+2au-a),\\
 \Delta(0,2,3,5)
  &=\frac{va^2(a-1)(b-1)}{(a-b)^2}(-Rb-a+2bu),\\
 \Delta(1,2,3,4)
  &=v(a-1)(b-1)(-Rb+a).
\end{split}
\]
All prefactors displayed here are nonzero.  Thus concyclicity gives
\[
 a(1-2u)+bR=0,\qquad
 a+b(R-2u)=0,\qquad
 a-bR=0.
\]
The third equation gives $a=bR$.  Substitution in the second gives
$2b(R-u)=0$, hence $R=u$.  The first equation then becomes
$2bu(1-u)=0$.  Since $b\ne0$ and $u=R>0$, one has $u=1$.
But $R=u=1$ and $R=u^2+v^2$ force $v=0$, contradicting
\eqref{eq:n6-pasch-coordinates}.  Thus $C(P)\geq8$.
\end{proof}

\subsection{Seven points and the planar Fano wall}

\begin{lemma}[Planar Fano wall]\label{lem:fano-seven-wall}
There do not exist seven distinct real points together with seven proper
four-point circles whose three-point complements form a Steiner triple
system on the seven points.
\end{lemma}

\begin{proof}
We first derive, rather than quote, the unique form of a Steiner triple
system on seven labels.  Name one triple $123$.  The two remaining triples
through $1$ partition $\{4,5,6,7\}$ into two pairs, so after relabelling they
are $145$ and $167$.  The triples through $2$ use neither pair $45$ nor pair
$67$; after possibly interchanging $6,7$, they are $246$ and $257$.  The
unused pairs on $\{4,5,6,7\}$ are then $47$ and $56$, forcing the final two
triples through $3$.  Thus the system is
\begin{equation}\label{eq:fano-triples}
 123,\quad145,\quad167,\quad246,\quad257,\quad347,
 \quad356.
\end{equation}

Invert at point $1$.  The four triples in \eqref{eq:fano-triples} avoiding
$1$ have complementary four-circles through $1$, which become the lines
\begin{equation}\label{eq:fano-inverted-lines}
 357,\qquad346,\qquad256,\qquad247.
\end{equation}
The first three form a nondegenerate triangle and the fourth is a
transversal.  Put
\[
 A=3,\quad B=5,\quad C=6,\qquad
 X=7\in AB,\quad Y=4\in AC,\quad Z=2\in BC.
\]
The complements of the three Fano triples through $1$ become the proper
circles
\begin{equation}\label{eq:fano-inverted-circles}
 ABYZ,\qquad ACXZ,\qquad BCXY.
\end{equation}

Use directed affine parameters
\[
 AB:A=0,B=1,X=x,\quad
 AC:A=0,C=1,Y=y,\quad
 BC:B=0,C=1,Z=z.
\]
All six quantities $x,y,z,1-x,1-y,1-z$ are nonzero.  Put
$a^2=|BC|^2$, $b^2=|CA|^2$, and $c^2=|AB|^2$.  Power of a point applied at
the opposite vertices to the three circles in
\eqref{eq:fano-inverted-circles} gives
\[
 b^2(1-y)=a^2(1-z),\qquad
 c^2(1-x)=a^2z,\qquad
 c^2x=b^2y.
\]
Multiplying the resulting ratios yields
\begin{equation}\label{eq:fano-power-product}
 \frac{x}{1-x}\frac{z}{1-z}\frac{1-y}{y}=+1.
\end{equation}
But $X,Y,Z$ are collinear by the fourth line in
\eqref{eq:fano-inverted-lines}, so directed Menelaus in triangle $ABC$
gives the same left-hand side equal to $-1$.  This contradiction proves the
lemma.
\end{proof}

\begin{proposition}\label{prop:n7}
One has $f(7)=11$.
\end{proposition}

\begin{proof}
Construction~\ref{con:seven-upper} gives $f(7)\leq11$.  Suppose that an
admissible seven-point set satisfies $C(P)\leq10$.  The two pencil bounds
give
\[
 \begin{array}{c|ccc}
 \text{circle size }m&4&5&6\\ \hline
 \text{lower bound}&7&15&13
 \end{array}
 \qquad
 \begin{array}{c|ccc}
 \text{line size }m&4&5&6\\ \hline
 \text{lower bound}&12&18&15.
 \end{array}
\]
The entry $m=4$ for circles need not be excluded; all larger circles and
all lines of size at least four are excluded.  Thus maximal blocks have size
three or four, and every four-block is a circle.

Let $L=L_3$, let $Q=C_4$, and let $R=C_3$.  Triple ownership gives
\begin{equation}\label{eq:n7-count}
 35=L+4Q+R,\qquad C(P)=Q+R=35-L-3Q.
\end{equation}
Since $21=L_2+3L$, Melchior gives $L_2\geq3$ and hence $L\leq6$.
The complements of two different four-circles intersect in at most one
point: the two circle supports intersect in at most two.  Therefore no pair
of labels occurs in two complementary triples, and $3Q\leq21$, so $Q\leq7$.

The assumption $C(P)\leq10$ and \eqref{eq:n7-count} imply
$L+3Q\geq25$.  If $Q\leq6$, the left side is at most $6+18=24$.
Consequently $Q=7$.  The seven complements contain
$7\binom32=21$ distinct pairs, hence every pair exactly once: they form a
Steiner triple system.  This contradicts \cref{lem:fano-seven-wall}.
\end{proof}

\subsection{Eight points}

\begin{lemma}[Even-arrangement slack]\label{lem:even-arrangement-slack}
Let $\mathcal A$ be a non-pencil arrangement of an even number of distinct
real projective lines.  Its Melchior slack
\[
 \kappa:=t_2-3-\sum_{r\geq4}(r-3)t_r
\]
cannot equal $1$.
\end{lemma}

\begin{proof}
Let $V,E,F$ be the numbers of vertices, edges, and faces of the projective
cell decomposition.  Euler's formula and the incidence count
$E=\sum_r r t_r$ give
\[
 \kappa=\sum_{\text{faces }G}(|G|-3).
\]
Because the number of lines is even, the sign of the product of homogeneous
linear equations for the lines is well-defined on the real projective plane
and two-colours the faces.  Every edge borders one face of each colour.
If $F_i$ is the number of faces of colour $i$, the sum of their boundary
lengths is $E$, and their excess is
\[
 \kappa_i=E-3F_i\geq0.
\]
Thus $\kappa_0\equiv\kappa_1\pmod3$ and
$\kappa=\kappa_0+\kappa_1\neq1$.
\end{proof}

\begin{proposition}\label{prop:n8}
One has $f(8)=17$.
\end{proposition}

\begin{proof}
Construction~\ref{con:eight-upper} gives the upper bound.  Suppose that an
admissible eight-point set has $C(P)\leq16$, and let $k$ be the maximum
number of selected points on a proper circle.  For $k=7,6,5$,
\cref{lem:rich-circle-pencil} gives, respectively,
\[
 C(P)\geq19,\qquad22,\qquad19.
\]
If $k=3$, every noncollinear triple has a different owner.  If a largest
line has $m\geq4$ points and $r=8-m$, the number of collinear triples is at
most
\[
 \binom m3+\binom r3+\binom r2\leq35.
\]
If all lines have size at most three, their triples use disjoint triples of
pairs and there are at most $\floor{28/3}=9$ of them.  In either event
$C(P)\geq21$.  The case $k=8$ is inadmissible, so only $k=4$ remains.

A line of size $m=5,6,7$ gives, by
\cref{lem:rich-line-pencil}, respectively $24,27,21$ circles.  Hence every
maximal generalized block has size three or four.  Put
\[
 b=B_4=C_4+L_4,\qquad
 L=L_3+L_4,\qquad h=L_4.
\]
Triple ownership, the summed pivot row, the original Melchior row, the
added-centre row, and the circle count are
\begin{align}
 B_3+4b&=56,                                         \label{eq:n8-triples}\\
 K:=\sum_{p\in P}\sigma_p&=144-12b\geq0,             \label{eq:n8-K}\\
 3L+4h&\leq25,                                       \label{eq:n8-global-line}\\
 9L+7h&\leq200-12b,                                  \label{eq:n8-added}\\
 C(P)&=56-3b-L.                                      \label{eq:n8-circle-count}
\end{align}
Indeed, \eqref{eq:n8-added} is
$9L_3+16L_4\leq56+K$ rewritten in terms of $L,h$.
Equations \eqref{eq:n8-K}, \eqref{eq:n8-global-line}, and
\eqref{eq:n8-circle-count} imply
\begin{equation}\label{eq:n8-b-range}
 b\in\{11,12\}.
\end{equation}

Suppose first that $b=11$.  Then $K=12$ and
\eqref{eq:n8-circle-count} gives $L\geq7$.  The two line inequalities force
\[
 L=7,\qquad h=0.
\]
At $p$, write $r_p=d_4(p)$.  Ownership of triples through $p$ and the
definition of $\sigma_p$ give
\[
 d_3(p)+3r_p=21,\qquad \sigma_p=18-3r_p.
\]
The local added-centre slack is
\[
 \kappa_p:=7+\sigma_p-3\ell_3(p).
\]
It is nonnegative and congruent to $1$ modulo $3$, hence at least one.
On the other hand,
\[
 \sum_{p\in P}\kappa_p=56+K-9L_3=56+12-63=5,
\]
which cannot be the sum of eight positive integers.

Now let $b=12$.  Then $K=0$, so $\sigma_p=0$ for every $p$, and
\begin{equation}\label{eq:n8-local-kappa}
 \kappa_p=7-3\ell_3(p)-4\ell_4(p)\geq0.
\end{equation}
The added-centre configuration has eight points and no line containing more
than four of them.  Its dual is therefore a non-pencil arrangement of eight
real projective lines, and $\kappa_p$ is its Melchior slack.  By
\cref{lem:even-arrangement-slack},
\begin{equation}\label{eq:n8-kappa-not-one}
 \kappa_p\neq1.
\end{equation}
Equations \eqref{eq:n8-local-kappa}--\eqref{eq:n8-kappa-not-one} imply
$\ell_3(p)\leq1$: if $\ell_4(p)=0$, the only additional nonnegative case is
$\ell_3(p)=2$, which gives slack one; if $\ell_4(p)\geq1$, then necessarily
$\ell_4(p)=1$ and already $\ell_3(p)\leq1$.  Hence
$3L_3\leq8$ and $L_3\leq2$.

Equation \eqref{eq:n8-circle-count} gives $L\geq4$.  Combining this with
\eqref{eq:n8-global-line}, \eqref{eq:n8-added}, and $L_3\leq2$ leaves
\begin{equation}\label{eq:n8-final-line-row}
 L_3=2,\qquad L_4=2,\qquad C(P)=16.
\end{equation}
Moreover,
\[
 \sum_p\kappa_p=56-9L_3-16L_4=6.
\]
The four possible local pairs and slacks are
\[
\begin{array}{c|cccc}
(\ell_3(p),\ell_4(p))&(0,0)&(1,0)&(0,1)&(1,1)\\ \hline
\kappa_p&7&4&3&0.
\end{array}
\]
Since $\sum_p\ell_4(p)=4L_4=8$, every point lies on one four-line; the slack
sum then forces two points of type $(0,1)$ and six of type $(1,1)$.  Thus the
two four-lines are disjoint and partition $P$.  Each of the two three-lines
contains three points, so two of them lie on the same four-line.  The
three-line and that four-line would then be the same geometric line, contrary
to maximality of their different supports.  This final contradiction proves
$C(P)\geq17$.
\end{proof}

\subsection{Nine points}

\begin{proposition}\label{prop:n9}
One has \(f(9)=25\).
\end{proposition}

\begin{proof}
Construction~\ref{con:circle-centre}, with four antipodal boundary pairs,
gives \(f(9)\leq1+\binom82-4=25\).  For the reverse inequality, suppose
that an admissible nine-point set \(P\) satisfies
\begin{equation}\label{eq:n9-counterassumption}
 C(P)\leq24,
\end{equation}
and let \(k\) be the largest number of selected points on a proper circle.
Then \(3\leq k\leq8\).

\smallskip
\noindent\emph{The branches \(k\geq6\).}
For \(k=8,7,6\), respectively, the rich-circle bound
\eqref{eq:rich-circle-pencil} gives
\[
 C(P)\geq25,\quad34,\quad28.
\]
Thus none of these branches is compatible with
\eqref{eq:n9-counterassumption}.

\smallskip
\noindent\emph{The branch \(k=5\).}
Fix a five-point circle \(\Gamma\) and put
\(X=P\setminus\Gamma\), so \(|X|=4\).  First,
\cref{lem:rich-line-pencil} excludes every line of size at least five:
for line sizes \(m=5,6,7,8\), its lower bound is, respectively,
\[
 28,\quad36,\quad39,\quad28.
\]
Hence all lines below have size at most four.

For a proper circle other than \(\Gamma\), let \(c_{r,t}\) count those
having \(r\) points on \(\Gamma\) and \(t\) points in \(X\).  Necessarily
\(r\leq2\).  Let \(\lambda_{r,t}\) denote the analogous number of rich
lines and set
\begin{equation}\label{eq:n9-R}
 R:=\sum_t t\lambda_{2,t}.
\end{equation}
At each outsider, the chords of \(\Gamma\) passing through it form a
matching on five vertices.  Therefore
\begin{equation}\label{eq:n9-R-cap}
 R\leq8.
\end{equation}

For \(x\in X\), let \(\mathcal S_x\) consist of the proper circles through
\(x\) and two points of \(\Gamma\).  Exactly the chord-lines through \(x\)
fail to yield such circles, whence
\[
 \sum_{x\in X}|\mathcal S_x|=40-R.
\]
For an outsider pair \(x,y\), the \(\Gamma\)-pairs represented in
\(\mathcal S_x\cap\mathcal S_y\) form a matching, so this intersection has
size at most two.  No circle contains two points of \(\Gamma\) and all four
outsiders, because \(k=5\).  Thus inclusion--exclusion stops after outsider
triples.  If \(c_2\) counts all circles with exactly two
\(\Gamma\)-points, then
\begin{equation}\label{eq:n9-four-outsider-union}
 c_2\geq40-R-12+c_{2,3}=28-R+c_{2,3}.
\end{equation}
Let \(c_{\leq1}\) count circles other than \(\Gamma\) with at most one
\(\Gamma\)-point.  Since \(C(P)=1+c_2+c_{\leq1}\),
\eqref{eq:n9-counterassumption} gives the four-outsider master inequality
\begin{equation}\label{eq:n9-four-outsider-master}
 R\geq5+c_{2,3}+c_{\leq1}.
\end{equation}

We now use only the four possible incidence types of \(X\).

\emph{No three outsiders are collinear and the four are not concyclic.}
The four outsider triples determine four distinct proper circles.  They
are counted by \(c_{2,3}+c_{\leq1}\), so
\eqref{eq:n9-four-outsider-master} gives \(R\geq9\), contradicting
\eqref{eq:n9-R-cap}.

\emph{All four outsiders are collinear.}
Put \(a=c_{1,2}\) and \(b=c_{2,2}\).  Counting the \(30\) triples with one
\(\Gamma\)-point and two outsiders gives
\begin{equation}\label{eq:n9-collinear-four-count}
 a+2b=30.
\end{equation}
There is no line contribution, since every outsider pair already lies on
the four-outsider line; nor can a proper circle contain three of its
collinear points.  The pair-intersection bound above gives \(b\leq12\),
and hence \(a\geq6\).  Moreover, the exact membership count for the
\(\mathcal S_x\) gives \(c_2=40-R-b\).  Therefore
\[
 C(P)=1+(40-R-b)+a
     =26-R+\frac32a\geq27,
\]
again a contradiction.

\emph{Exactly three outsiders are collinear.}
Let \(\lambda\) be their line and let \(s\in\{0,1\}\) record whether
\(\lambda\) contains a point of \(\Gamma\).  Put
\[
 a:=c_{1,3},\qquad h:=c_{2,3},\qquad
 \ell:=\lambda_{1,2}.
\]
The other three outsider triples determine three distinct proper circles.
Thus \(h+c_{\leq1}\geq3\).  Combining this with
\eqref{eq:n9-R-cap}, \eqref{eq:n9-four-outsider-union}, and
\eqref{eq:n9-four-outsider-master} first gives
\(R=8\) and \(h+c_{\leq1}=3\).  If
\[
 I:=\sum_{\{x,y\}\subset X}|\mathcal S_x\cap\mathcal S_y|
   =c_{2,2}+3h,
\]
then \(I\leq12\), while exact inclusion--exclusion gives
\[
 C(P)=1+40-R-I+h+c_{\leq1}=36-I.
\]
Thus \eqref{eq:n9-counterassumption} forces \(I=12\) and equality
throughout:
\begin{equation}\label{eq:n9-three-collinear-equality}
 R=8,\qquad C(P)=24,\qquad
 h+c_{\leq1}=3,\qquad c_{2,2}+3h=12.
\end{equation}
In particular those three outsider-triple circles exhaust
\(c_{\leq1}+h\), so \(c_{1,2}=0\).  The last equality says that each
outsider pair lies on exactly two circles with two \(\Gamma\)-points, and
their two \(\Gamma\)-pairs are disjoint.  There can be no line of type
\((2,2)\): its \(\Gamma\)-pair would have to be disjoint from both circle
pairs, which is impossible on five points.

Counting again the triples with one \(\Gamma\)-point and two outsiders,
and then applying the global line row, gives
\begin{align}
 \ell+3a+3s&=6,                         \label{eq:n9-three-collinear-triples}\\
 3\ell+4s&\leq6.                        \label{eq:n9-three-collinear-lines}
\end{align}
Indeed, the term already covered by the two-\(\Gamma\) circles is
\(2(c_{2,2}+3h)=24\); while the rich lines are the eight chord incidences,
\(\lambda\), and the \(\ell\) lines of type \((1,2)\).
Equations \eqref{eq:n9-three-collinear-triples}--%
\eqref{eq:n9-three-collinear-lines} yield
\[
 \ell=0,\qquad a+s=2.
\]
For each outsider pair, the two disjoint \(\Gamma\)-pairs in
\eqref{eq:n9-three-collinear-equality} leave one \(\Gamma\)-point.  Its
triple with the outsider pair must be owned by one of the \(a\)
three-outsider circles or, when \(s=1\), by \(\lambda\).
The pair-shadows of two distinct three-subsets of a four-set cover only
five of its six pairs.  One pair is therefore left without an owner, a
contradiction.

\emph{The four outsiders are concyclic.}
Let \(\Omega\) be their circle.  It contains \(\delta\in\{0,1\}\) selected
points of \(\Gamma\).  Put
\[
 a:=c_{1,2},\quad b:=c_{2,2},\quad
 u:=\lambda_{2,1},\quad v:=\lambda_{2,2},\quad
 \ell:=\lambda_{1,2}.
\]
Apart from \(\Gamma,\Omega\), every rich circle or line has type
\((2,1),(1,2)\), or \((2,2)\): a distinct carrier meets each of
\(\Gamma,\Omega\) in at most two points.  Consequently
\begin{align}
 R&=u+2v,                                                   \label{eq:n9-concyclic-R}\\
 C(P)&=42-R-b+a.                                           \label{eq:n9-concyclic-C}
\end{align}
Write \(b=12-e\), where \(e\geq0\).  From
\eqref{eq:n9-counterassumption}, \eqref{eq:n9-R-cap}, and
\eqref{eq:n9-concyclic-C},
\begin{equation}\label{eq:n9-concyclic-master}
 R\geq6+e+a,\qquad e+a\leq2.
\end{equation}
The same \(30\)-triple count and the global line row now read
\begin{align}
 \ell+2v&=6-a+2e-6\delta,                 \label{eq:n9-concyclic-triples}\\
 3R+3\ell+v&\leq33.                       \label{eq:n9-concyclic-lines}
\end{align}

If \(\delta=0\), all \(b+v\) generalized blocks of type \((2,2)\)
are cross-blocks between five points of \(\Gamma\) and four points of
\(\Omega\).  By \cref{lem:radical-axis-cross-block},
\(b+v\leq12\), hence \(v\leq e\).  But
\eqref{eq:n9-concyclic-master}--\eqref{eq:n9-concyclic-lines} give
\[
 33\geq3R+3\ell+v\geq36+9e-5v,
\]
or \(5v\geq3+9e\), impossible when \(v\leq e\).

If \(\delta=1\), remove the common selected point of \(\Gamma\) and
\(\Omega\).  A type-\((2,2)\) block cannot use that point: a circle would
share three points with \(\Omega\), while a line would contain three
collinear points of \(\Omega\).  The \(4+4\) case of
\cref{lem:radical-axis-cross-block} therefore gives
\[
 b+v\leq10,\qquad v\leq e-2.
\]
Together with \eqref{eq:n9-concyclic-master}, this forces
\[
 e=2,\qquad a=v=0,\qquad R=8.
\]
Equation~\eqref{eq:n9-concyclic-triples} gives \(\ell=4\), and
\eqref{eq:n9-concyclic-lines} becomes \(36\leq33\), the final
contradiction in the \(k=5\) branch.

\smallskip
\noindent\emph{The branches \(k\leq4\).}
The same rich-line calculation excludes lines of size at least five.
All maximal generalized blocks therefore have size three or four.  Put
\[
 B:=B_4=C_4+L_4,\qquad L:=L_3+L_4.
\]
Triple ownership and the circle objective give
\begin{equation}\label{eq:n9-four-block-count}
 C(P)=84-3B-L.
\end{equation}
The global line row is \(3L_3+7L_4\leq33\), so \(L\leq11\).
Under \eqref{eq:n9-counterassumption},
\eqref{eq:n9-four-block-count} then forces \(B\geq17\).

Fix \(p\in P\) and invert the other eight points at \(p\).  They are
noncollinear by \cref{lem:inversion-dictionary}, and no four are collinear:
such a line would invert to a five-point proper circle or a five-point
line through \(p\).  A three-point image line is exactly a generalized
four-block of \(P\) through \(p\).  Hence
\cref{lem:orchard-eight} gives \(d_4(p)\leq7\).  Summing over all pivots,
\[
 4B=\sum_{p\in P}d_4(p)\leq9\cdot7=63,
\]
so \(B\leq15\), contrary to \(B\geq17\).

Every possible value of \(k\) has now been excluded under
\eqref{eq:n9-counterassumption}.  Thus \(C(P)\geq25\), completing the
proof.
\end{proof}

\begin{theorem}[The small values]\label{thm:small-values}
For the V1 problem,
\[
 \bigl(f(4),f(5),f(6),f(7),f(8),f(9)\bigr)
   =(3,5,8,11,17,25).
\]
\end{theorem}

\begin{proof}
Combine \cref{prop:n4-n5,prop:n6,prop:n7,prop:n8,prop:n9}.
\end{proof}
